% 第17批（自测题6）：选择6.1.1—6.1.10；原书印刷163—164/PDF171—172。
\subsection{自测题6：电磁场基础选择题}

\sourcetest{6.1.1}{163}{171}
\problemtext $x$ 轴上左半无限长线电荷密度 $+\lambda$、右半无限长密度 $-\lambda$，求 $P(0,a)$ 的电场。
\workings 由库仑场沿两段线分别积分，竖直分量抵消，而水平分量同向。每段贡献 $E_x=\lambda/(4\pi\varepsilon_0a)$，故
\[\boxed{\bm E(0,a)=\frac{\lambda}{2\pi\varepsilon_0a}\,\bm i}.\]
原选择题选 \textbf{B}。
\crosscheck 正电左段排斥、负电右段吸引，二者都指向 $+x$；若只把两段净线电荷“相加为0”，会错误地得出零场。

\sourcetest{6.1.2}{163}{171}
\problemtext 正点电荷 $+q$ 到 $P$ 的距离为 $a$、到 $M$ 的距离为 $2a$，把 $P$ 选作电势零点，求 $M$ 电势。
\workings 以无穷远为零时 $\varphi(r)=q/(4\pi\varepsilon_0r)$；变更电势参考值要整体减去 $\varphi(P)$，故
\[\boxed{\varphi_M=\frac{q}{4\pi\varepsilon_0}\left(\frac1{2a}-\frac1a\right)=-\frac{q}{8\pi\varepsilon_0 a}}.\]
选 \textbf{D}。
\crosscheck 电势差不依参考零点；若直接写正值 $q/(8\pi\varepsilon_0a)$，忽略了题目重新设定的零势点。

\sourcetest{6.1.3}{163}{171}
\problemtext 同心薄球壳，内壳半径 $R_1$ 带电 $Q$，外壳半径 $R_2$ 最初中性但接地；求两壳间 $R_1<r<R_2$ 的场强与电势。
\workings 对高斯球，$\bm E=\frac{Q}{4\pi\varepsilon_0r^2}\hat{\bm r}$。外壳接地意味着 $V(R_2)=0$，对径向电场积分：
\[\boxed{V(r)=\frac{Q}{4\pi\varepsilon_0}\left(\frac1r-\frac1{R_2}\right)},\quad \boxed{E(r)=\frac{|Q|}{4\pi\varepsilon_0r^2}}.\]
与选项 \textbf{C} 相符。
\crosscheck 电势参考是接地外壳，不是无穷远；场强大小为正，方向由 $Q$ 符号决定。

\sourcetest{6.1.4}{163}{171}
\problemtext 等边三角形边长 $a$，三顶点分别放电荷 $q,2q,3q$，将另一电荷 $Q$ 从无穷远缓慢移到中心，求外力功。
\workings 三顶点到重心的距离同为 $a/\sqrt3$，因此中心电势 $V_O=\frac{1}{4\pi\varepsilon_0}\frac{(q+2q+3q)\sqrt3}{a}$。准静态移入外力功 $\Delta U=QV_O$，即
\[\boxed{W_{\rm ext}=\frac{6\sqrt3\,qQ}{4\pi\varepsilon_0a}}.\]
选 \textbf{C}。
\crosscheck 应当相加的是标量电势，不是中心处三个矢量电场。

\sourcetest{6.1.5}{163}{171}
\problemtext 质子在匀强磁场 $B=0.3$T 中垂直于磁场作半径 $R=10$cm 的圆周运动，估计动能数量级（质量 $1.67\times10^{-27}$kg，电荷 $1.6\times10^{-19}$C）。
\workings 洛伦兹力充当向心力：$qvB=mv^2/R$，故 $p=mv=qBR$、$K=q^2B^2R^2/(2m)$。代入
\[\boxed{K\approx6.90\times10^{-15}\ {\rm J}\approx0.0431\ {\rm MeV}}.\]
在题给 $0.01$、$0.1$、$1$、$10$、$100$MeV 中数量级最接近 \textbf{B}（$0.1$MeV）。
\crosscheck 题问的是数量级，而不是精确数值；$v=qBR/m\simeq2.87\times10^6$ m/s，使用非相对论动能成立。

\sourcetest{6.1.6}{164}{172}
\problemtext $N$ 匝边长为 $a$ 的正三角形通电线圈位于匀强磁场中，线圈法线与 $\bm B$ 同向，求磁力矩大小。
\workings 任意平面线圈磁矩 $\bm m=NI S\,\hat{\bm n}$，其中 $S=\sqrt3 a^2/4$。磁力矩 $\bm\tau=\bm m\times\bm B$，因 $\hat{\bm n}\parallel\bm B$，
\[\boxed{|\bm\tau|=NISB\sin0=0}.\]
选 \textbf{D}。
\crosscheck 有磁通和磁矩不代表有磁力矩；只有磁矩与磁场夹角非零才有转矩。

\sourcetest{6.1.7}{164}{172}
\problemtext 无限长薄平电流带宽度 $a$，总电流 $I$ 沿带长方向均匀分布，带外离近边缘 $b$ 的点 $P$ 磁场多大？
\workings 把电流带分解成宽 $dx$ 的无限长线电流 $dI=(I/a)\,dx$。同侧场方向一致，取从近边界量起的坐标 $x\in[0,a]$：
\[B(P)=\frac{\mu_0 I}{2\pi a}\int_0^a\frac{dx}{b+x}
=\boxed{\frac{\mu_0 I}{2\pi a}\ln\frac{a+b}{b}}.\]
选 \textbf{B}。
\crosscheck 必须按每条细线到 $P$ 的距离 $b+x$ 积分；直接把全部电流集中在中心只给近似。

\sourcetest{6.1.8}{164}{172}
\problemtext 原一匝圆线圈半径 $R$、电流 $I$，将原导线在总长度不变下绕成两匝同向圆线圈，并通以同样的 $I$；求中心磁场与磁矩变化倍数。
\workings 原长 $2\pi R$，新两匝各半径 $R'=R/2$。中心磁场 $B'=\mu_0NI/(2R')=\mu_0(2I)/R=4B$；磁矩大小 $m'=NI\pi R'^2=2I\pi R^2/4=m/2$。故
\[\boxed{\frac{B'}B=4,\qquad \frac{m'}m=\frac12}.\]
选 \textbf{B}。
\crosscheck 中心磁场对匝数和半径均敏感；磁矩与匝数乘面积成正比。

\sourcetest{6.1.9}{164}{172}
\problemtext 导体棒 $AB$ 绕棒上的 $C$ 点、以垂直棒并与匀强磁场平行的轴匀速旋转；$BC=L/3$，比较端点 $A,B$ 的电势。
\workings 取 $\bm\omega$ 与 $\bm B$ 同向，棒上电荷感受 $\bm v\times\bm B$ 沿远离旋转轴的径向。端点相对轴的距离 $CA=2L/3$、$CB=L/3$，所以
\[\boxed{V_A-V_B=\tfrac12B\omega(CA^2-CB^2)=\frac{B\omega L^2}{6}>0}.\]
按原图转向选 \textbf{A}。
\crosscheck 电势并非单由端点速度线性比较，应按 $\omega Br\,dr$ 积分；若反转角速度与磁场相对方向，电势高低反转。

\sourcetest{6.1.10}{164}{172}
\problemtext 导体杆在垂直纸面向里的匀强磁场中沿轨道向右匀加速运动，外接铁芯绕组和电容器；判断达到稳定运动阶段后电容器 $M$ 极板的电荷符号和随时间变化。
\workings 磁通增加时感应电流方向由楞次定律确定：回路倾向产生纸面向外的磁场。若电容器经铁芯耦合，其极板电势要继续使用绕组的实际接线次序和感应电动势的符号判定。只根据杆运动方向可确认其来源电动势 $|\mathcal E|=B\ell v$ 随 $v$ 增大，却\emph{不能}不看绕组极性而确定 $M$ 正负。这里保留 \textbf{源图接线二校}，不擅定 A/B/C/D。
\crosscheck 必须逐匝确认绕组绕向；电容器是否持续充电还取决于铁芯磁导率及电路稳态近似，不能把理想磁耦合直接等同导线相接。
